% RQ3 Introduction: Proposal Pipeline Effects

\subsection{Motivation}

Every DAO proposal follows a pipeline: a sequence of stages with gates, timers, and decision points. This pipeline mediates between the raw demand for governance decisions and the organization's capacity to process them thoughtfully.

Pipeline design involves fundamental tradeoffs:
\begin{itemize}
    \item \textbf{Speed vs. deliberation}: Faster decisions reduce opportunity for input and analysis
    \item \textbf{Filtering vs. inclusion}: Strict gates may exclude worthy proposals
    \item \textbf{Flexibility vs. predictability}: Expedited paths may be abused or create precedent confusion
\end{itemize}

Real DAOs have evolved diverse pipeline approaches. Compound uses a linear flow (proposal $\rightarrow$ voting $\rightarrow$ timelock). Optimism employs multi-stage review. Many DAOs use off-chain ``temp-checks'' via Snapshot before formal on-chain governance.

Despite this diversity, little systematic analysis exists of how pipeline parameters affect outcomes. Practitioners rely on intuition or copy existing patterns without understanding tradeoffs.

\subsection{Research Question}

This paper investigates:

\begin{quote}
\textbf{RQ3:} How do proposal pipeline settings (temp-checks, fast-tracks, expiry windows) affect throughput and decision quality?
\end{quote}

Specifically, we examine:
\begin{enumerate}
    \item How do temp-check thresholds affect the volume and quality of proposals reaching formal voting?
    \item When do fast-track provisions improve governance efficiency vs. create risks?
    \item How do expiry windows interact with deliberation time and proposal abandonment?
    \item What pipeline configurations optimize for different DAO objectives?
\end{enumerate}

\subsection{Pipeline Mechanisms}

We study three pipeline mechanisms:

\subsubsection{Temp-Checks}

Preliminary signal votes (often off-chain) that filter proposals before formal governance:
\begin{equation}
\text{advance}(p) = \mathbf{1}[\text{temp-check support}(p) > \tau_{\text{temp}}]
\end{equation}

\subsubsection{Fast-Tracks}

Expedited paths for proposals meeting certain criteria:
\begin{itemize}
    \item High initial support (``obvious'' consensus)
    \item Emergency/time-sensitive matters
    \item Routine operational decisions
\end{itemize}

\subsubsection{Expiry Windows}

Maximum time allowed for proposal completion:
\begin{equation}
\text{expired}(p) = \mathbf{1}[\text{age}(p) > \tau_{\text{expiry}}]
\end{equation}

Expired proposals fail without resolution.

\subsection{Contributions}

This paper contributes:

\begin{enumerate}
    \item \textbf{Systematic analysis} of pipeline mechanism effects on governance outcomes
    \item \textbf{Throughput-quality metrics} for evaluating pipeline configurations
    \item \textbf{Parameter guidelines} for temp-checks, fast-tracks, and expiry windows
    \item \textbf{Configuration recommendations} for different DAO operational contexts
\end{enumerate}

\subsection{Paper Organization}

Section~\ref{sec:background} reviews proposal pipeline practices in deployed DAOs. Section~\ref{sec:theory} develops formal models of pipeline dynamics. Section~\ref{sec:architecture} describes simulation implementation. Section~\ref{sec:methodology} details experimental design. Section~\ref{sec:results} presents findings. Section~\ref{sec:discussion} interprets results for practitioners. Section~\ref{sec:conclusion} concludes.
